%=======================================================================
% CSE465 Final Report - OWMTL
% FORMAT SKELETON ONLY. See paper_requirement.md before writing content.
%
% Supervisor hard rules (slides 2 and 4 of 4X5_Final_Report_261):
%   - IEEE single-column journal format is MANDATORY (line below).
%   - Compile on Overleaf. Submit the Overleaf link with EDIT access.
%   - Minimum 6 pages.
%   - LaTeX Error => not allowed. LaTeX Warning => allowed.
%=======================================================================
\documentclass[journal,onecolumn]{IEEEtran}

\usepackage{cite}                 % IEEE-style [1], [2]-[4] collapsing
\usepackage{amsmath,amssymb}
\usepackage{graphicx}
\usepackage{booktabs}
\usepackage{multirow}
\usepackage{url}
\usepackage[hidelinks]{hyperref}  % load last

\graphicspath{{figures/}}

% Placeholder box so the skeleton compiles before real figures exist.
% DELETE this macro and every \figstub call once figures are dropped in.
\newcommand{\figstub}[1]{%
  \fbox{\begin{minipage}[c][4.5cm][c]{0.92\linewidth}\centering\textit{#1}\end{minipage}}}

\begin{document}

%=======================================================================
\title{TODO: Paper title\\
{\normalsize\itshape (working title is in PAPER\_OUTLINE.md)}}

\author{Asif~TODO,
        Barshon~Basak,
        Farhana~TODO,
        and~Sami~TODO%
\thanks{Department of Electrical and Computer Engineering, North South
University, Dhaka, Bangladesh. Course: CSE465, Section TODO.
Supervisor: TODO. E-mail: TODO.}}

\maketitle

%=======================================================================
% ABSTRACT -- 200-300 words. No citations, no abbreviations, no symbols,
% no equations. Cover: purpose -> dataset -> preprocessing -> models ->
% results -> ablation/XAI -> conclusion.
%=======================================================================
\begin{abstract}
TODO.
\end{abstract}

% Keywords: main topics, ALPHABETICAL ORDER.
\begin{IEEEkeywords}
TODO, TODO, TODO, TODO, TODO.
\end{IEEEkeywords}

%=======================================================================
\section{Introduction}
\label{sec:intro}
% P1  -- motivation: why this project was chosen.
% P2+ -- literature review. At least one paper reviewed per group member
%        (4 total), journal preferred over conference. Cite as \cite{ref1}.
% P_n -- END of the review: gaps, limitations, weaknesses of prior work.
% next 3-4 paragraphs -- YOUR work.
% Last paragraph -- roadmap of the remaining sections.

TODO: motivation \cite{ref1}.

The remainder of this paper is organised as follows. Section~\ref{sec:proposed}
describes the proposed system. Section~\ref{sec:results} presents the results
and discussion. Section~\ref{sec:conclusion} concludes the paper.

%=======================================================================
\section{Proposed System}
\label{sec:proposed}
% NO RESULTS IN THIS SECTION. Methodology only.
TODO: one-paragraph overview, then the system flowchart.

\begin{figure}[!t]
\centering
\figstub{Fig. 1 -- flowchart of the complete system}
% \includegraphics[width=0.95\linewidth]{system_flowchart.pdf}
\caption{Flowchart of the complete proposed system.}
\label{fig:flowchart}
\end{figure}

The end-to-end pipeline is summarised in Fig.~\ref{fig:flowchart}.

\subsection{Dataset}
\label{subsec:dataset}
% Cite the dataset itself.
TODO: source, size, classes, split policy \cite{ref2}.

\begin{table}[!t]
\centering
\caption{Class distribution before and after data augmentation.}
\label{tab:dataset}
\begin{tabular}{lrrrr}
\toprule
Class & Open-source & Collected & Total & Augmented total \\
\midrule
TODO & 0 & 0 & 0 & 0 \\
TODO & 0 & 0 & 0 & 0 \\
\bottomrule
\end{tabular}
\end{table}

Table~\ref{tab:dataset} reports the class distribution.

\subsection{Dataset Preprocessing}
\label{subsec:preprocessing}
% ALL preprocessing techniques must be applied and described here.
% Training-set data augmentation belongs here too.
TODO. Every quantity introduced here gets an equation, e.g.

\begin{equation}
\label{eq:example}
x_{\mathrm{norm}} = \frac{x - \mu}{\sigma},
\end{equation}

\noindent where (\ref{eq:example}) is applied per TODO.

\subsection{Applied Models}
\label{subsec:models}
% One pre-trained model per group member (4), PLUS at least 4 transformer
% models (ViT / Swin / DINO / PVT / CvT / DeiT ...). A custom model or
% algorithm gets its own figure and/or algorithm block here.
TODO.

%=======================================================================
\section{Results and Discussion}
\label{sec:results}
% The most important section. Every number lives here.
TODO.

\begin{table}[!t]
\centering
\caption{Classification performance of all applied models.}
\label{tab:performance}
\begin{tabular}{lrrrr}
\toprule
Model & Accuracy & Precision & Recall & F1-score \\
\midrule
TODO & 0.0000 & 0.0000 & 0.0000 & 0.0000 \\
TODO & 0.0000 & 0.0000 & 0.0000 & 0.0000 \\
\bottomrule
\end{tabular}
\end{table}

\begin{table}[!t]
\centering
\caption{Model size, parameter count, and training time. Training was
performed on TODO (GPU / runtime).}
\label{tab:complexity}
\begin{tabular}{lrrrr}
\toprule
Model & Epochs & Time (min) & Params & Size (MB) \\
\midrule
TODO & 0 & 0.00 & 0 & 0.00 \\
TODO & 0 & 0.00 & 0 & 0.00 \\
\bottomrule
\end{tabular}
\end{table}

Tables~\ref{tab:performance} and~\ref{tab:complexity} summarise TODO.

% --- Figures below are required for the BEST model only ---
\begin{figure}[!t]
\centering
\figstub{Fig. 2 -- normalized confusion matrix, best model}
\caption{Normalized confusion matrix of the best-performing model.}
\label{fig:cm}
\end{figure}

\begin{figure}[!t]
\centering
\figstub{Fig. 3 -- training/validation loss and accuracy vs. epochs}
\caption{Training and validation loss and accuracy versus epochs for the
best-performing model.}
\label{fig:curves}
\end{figure}

\subsection{Explainable AI}
\label{subsec:xai}
% LIME / SHAP / Grad-CAM on the BEST model. Interpret it, do not just show it.
\begin{figure}[!t]
\centering
\figstub{Fig. 4 -- XAI visualisation}
\caption{TODO XAI visualisation for the best-performing model.}
\label{fig:xai}
\end{figure}

Fig.~\ref{fig:xai} shows TODO.

\subsection{Ablation Study}
\label{subsec:ablation}
% One row per metric, with a "possible reason" column explaining the delta.
\begin{table}[!t]
\centering
\caption{Ablation study on the best-performing model.}
\label{tab:ablation}
\begin{tabular}{lrrl}
\toprule
Metric & With TODO & Without TODO & Possible reason \\
\midrule
Accuracy & 0.0000 & 0.0000 & TODO \\
F1-score & 0.0000 & 0.0000 & TODO \\
\bottomrule
\end{tabular}
\end{table}

Table~\ref{tab:ablation} reports TODO.

\subsection{Comparison with Existing Works}
\label{subsec:comparison}
% MANDATORY, and it goes at the END of Results and Discussion.
% The last row is always "This work".
\begin{table}[!t]
\centering
\caption{Performance comparison of the proposed system with existing works.}
\label{tab:comparison}
\begin{tabular}{llrrr}
\toprule
Author & Dataset & Model & Accuracy & F1-score \\
\midrule
\cite{ref1} & TODO & TODO & 0.0000 & 0.0000 \\
\cite{ref2} & TODO & TODO & 0.0000 & 0.0000 \\
\midrule
This work & TODO & TODO & 0.0000 & 0.0000 \\
\bottomrule
\end{tabular}
\end{table}

Table~\ref{tab:comparison} compares TODO.

%=======================================================================
\section{Conclusions}
\label{sec:conclusion}
% Close the paper, then TWO or THREE future works.
TODO.

%=======================================================================
\bibliographystyle{IEEEtran}
\bibliography{references}

\end{document}
